% UTP FYP technical paper. Layout comes from reference/UTP-FYP-Technical-Paper-Template.tex
% (UTP FYP Guidelines v7.0, Sections 5.9 and 6.2, Appendix 10-1). Limit: 6 pages, one column,
% IEEE references, and every reference listed must be cited in the text (and vice versa).
% Build from this folder with `latexmk`. The engine must be XeLaTeX; the PDF goes to build/.

\documentclass[a4paper,11pt]{article}

% ------------------------------------------------------------
% FONT SETUP
% ------------------------------------------------------------
% Guideline 6.2(c): Times New Roman 11 (or Arial 10).
\usepackage{iftex}
\ifPDFTeX
  \usepackage[T1]{fontenc}
  \usepackage[utf8]{inputenc}
  \usepackage{newtxtext,newtxmath}
\else
  \usepackage{fontspec}
  \IfFontExistsTF{Times New Roman}{%
    \setmainfont{Times New Roman}%
  }{%
    \IfFontExistsTF{Tinos}{%
      \setmainfont{Tinos}%
    }{%
      \setmainfont{Nimbus Roman}%
    }%
  }
  \usepackage{unicode-math}
  \IfFontExistsTF{Times New Roman}{%
    \setmathfont{XITSMath-Regular.otf}%
  }{}
\fi

% ------------------------------------------------------------
% CORE PACKAGES
% ------------------------------------------------------------
% Guideline 6.2(d): 1 in (2.5 cm) margins on all sides.
\usepackage[a4paper,margin=2.5cm]{geometry}
\usepackage{setspace}
\usepackage{titlesec}
\usepackage{caption}
\usepackage{graphicx}
\graphicspath{{figures/}{../results/}}
\usepackage{booktabs}
\usepackage{array}
\usepackage{tabularx}
\usepackage{float}
\usepackage{amsmath}
\usepackage{enumitem}
\usepackage{etoolbox}
\usepackage{csquotes}
\usepackage[hidelinks]{hyperref}

% ------------------------------------------------------------
% IEEE REFERENCES
% ------------------------------------------------------------
% Reads references.bib beside this .tex file.
% Use \cite{key} in the text, which gives [1].
\usepackage[style=ieee,backend=biber]{biblatex}
\addbibresource{references.bib}
\AtBeginBibliography{\sloppy}

% ------------------------------------------------------------
% PAPER METADATA
% ------------------------------------------------------------
% TODO: Replace these fields. The title block reads from them.
% The title should be less than 15 words.
\newcommand{\PaperTitle}{FinBERT News Sentiment for Stock Price Movement Prediction in the US Market}
\newcommand{\WriterName}{Mas Hairul Adham Bin Alwee}
\newcommand{\CoWriterName}{Nurul Aida Binti Osman}
% TODO: Confirm the department name.
\newcommand{\AffiliationOne}{Computer and Information Sciences Department, Universiti Teknologi PETRONAS, 32610 Bandar Seri Iskandar, Perak, Malaysia}
\newcommand{\CorrespondingEmail}{mas\_22005746@utp.edu.my}

% ------------------------------------------------------------
% BODY SPACING AND PARAGRAPH STYLE
% ------------------------------------------------------------
% 1.5 spacing, justified, flush-left paragraphs with one blank line between.
\onehalfspacing
\setlength{\parindent}{0pt}
\setlength{\parskip}{11pt plus 2pt minus 1pt}
\raggedbottom

% Tables and figures: single spacing.
\AtBeginEnvironment{table}{\singlespacing\setlength{\parskip}{0pt}}
\AtBeginEnvironment{figure}{\singlespacing\setlength{\parskip}{0pt}}

\setlist{topsep=0pt, parsep=0pt, itemsep=0pt, partopsep=0pt}

% ------------------------------------------------------------
% PAGINATION
% ------------------------------------------------------------
\pagestyle{plain}

% ------------------------------------------------------------
% HEADING LEVELS
% ------------------------------------------------------------
% Appendix 10-1 uses unnumbered headings, at most three levels:
%   \section       HEADING (bold, uppercase)
%   \subsection    Subheading (bold)
%   \subsubsection Subheading (bold italic)
\setcounter{secnumdepth}{0}

\titleformat{\section}
  {\normalfont\bfseries\normalsize}
  {}
  {0pt}
  {\MakeUppercase}
\titlespacing*{\section}{0pt}{12pt}{0pt}

\titleformat{\subsection}
  {\normalfont\bfseries\normalsize}
  {}
  {0pt}
  {}
\titlespacing*{\subsection}{0pt}{12pt}{0pt}

\titleformat{\subsubsection}
  {\normalfont\bfseries\itshape\normalsize}
  {}
  {0pt}
  {}
\titlespacing*{\subsubsection}{0pt}{12pt}{0pt}

% ------------------------------------------------------------
% FIGURES AND TABLES
% ------------------------------------------------------------
% "Table 1 Title" above tables (11 pt) and "Figure 1 Title" below
% figures (10 pt legend), bold label, no colon.
\DeclareCaptionFont{tenpt}{\fontsize{10pt}{12pt}\selectfont}
\captionsetup{labelfont=bf,labelsep=space,justification=centering,singlelinecheck=true}
\captionsetup[table]{position=above,skip=4pt}
\captionsetup[figure]{position=below,font=tenpt,skip=6pt}

% Optional thin border around an inserted image:
%   \FramedGraphic{width=0.75\textwidth}{assets/my-image.png}
\newcommand{\FramedGraphic}[2]{%
  \setlength{\fboxsep}{0pt}%
  \setlength{\fboxrule}{0.5pt}%
  \fbox{\includegraphics[#1]{#2}}%
}

% Table column type: left-aligned, stretches to fill the table width.
\newcolumntype{L}{>{\raggedright\arraybackslash}X}

% ============================================================
% DOCUMENT START
% ============================================================
\begin{document}

% ------------------------------------------------------------
% TITLE BLOCK
% ------------------------------------------------------------
% For more authors/affiliations, add \textsuperscript{2}, {3}, ... after
% each name and one \textsuperscript{n}Affiliation line per number.
\begin{center}
\singlespacing
{\bfseries \PaperTitle\par}
\vspace{6pt}
{\bfseries \WriterName\textsuperscript{1}, \CoWriterName\textsuperscript{1}\par}
\vspace{6pt}
\textsuperscript{1}\AffiliationOne\par
\vspace{6pt}
Corresponding email: \CorrespondingEmail\par
\end{center}


% ------------------------------------------------------------
% ABSTRACT (single spacing, 11 pt, no citations, < 300 words)
% ------------------------------------------------------------
\vspace{12pt}
{\bfseries Abstract\par}
{\singlespacing
Financial news sentiment is widely used to supplement stock price prediction, but its contribution is often reported without a matched price-only baseline or across a single news source. This study measured the contribution of FinBERT sentiment to stock movement prediction, using 34,810 FNSPID headlines from 2018 to 2023 for five large-capitalisation US stocks. XGBoost and LSTM models were trained with and without sentiment features under identical settings, and were tested by walk-forward evaluation on 2021--2023 at horizons of one, three and five trading days. Six other non-agentic sentiment models were also compared with FinBERT. Differences were tested with a paired block bootstrap and corrected for multiple comparisons. Sentiment did not significantly improve XGBoost or a class-weighted LSTM, and a significant gain from a plain LSTM was shown to be an artifact of class imbalance. All models remained close to chance. As publication times were unavailable for 99.8\% of the headlines, the experiments were repeated with sentiment lagged by one day, and the conclusion was unchanged. No sentiment model was significantly better than another. Sentiment was correlated with the same day's return for all five stocks, but much less with later returns, and a trading rule based on it lost money after costs. Agreement between sentiment and price direction varied across outlets when the publication day was included, but fell to about 50\% on later returns. The findings are consistent with the efficient market hypothesis, and suggest that headline sentiment reflects price moves as they occur and provides no reliable information about later moves.
\par}

\vspace{6pt}
\textbf{\textit{Keywords}}\quad \textit{ablation study; XGBoost; long short-term memory; walk-forward evaluation; FNSPID; multiple comparisons}

\section{Introduction}
Stock prices respond to both quantitative signals, such as historical prices and volumes, and qualitative signals, such as news. FinBERT, a BERT model adapted to financial text, has become a standard tool for extracting sentiment from news \cite{araci2019finbert}. It has been combined with sequential and tree-based models for stock prediction \cite{gu2024predicting,mathematics2025stock}, but two gaps recur in this work. Many studies evaluate a single news source, although FinBERT accuracy has been reported to vary from 32\% to 76\% across five outlets \cite{liu2025enhancing}. Many also report a sentiment-augmented model without a price-only model evaluated under the same conditions, so the contribution of sentiment cannot be isolated.

This study addressed both gaps. An ablation design was used to measure the contribution of sentiment over a matched price-only baseline at three horizons, and FinBERT sentiment was examined across news outlets. FinBERT was also compared with six other non-agentic sentiment models, on labelled text and on the prediction task.

\section{Literature Review}
The efficient market hypothesis holds that public information is rapidly incorporated into prices, which would limit the value of news sentiment \cite{fama1970efficient}. FinBERT was developed because general sentiment models perform poorly on financial language \cite{araci2019finbert,yang2020finbert}, and it is commonly validated on the Financial PhraseBank \cite{malo2014good}. A FinBERT-LSTM has been reported to improve on price-only baselines \cite{gu2024predicting}, and the authors of the FNSPID dataset reported a modest gain from sentiment \cite{dong2024fnspid}. In large finance studies, the price response to news has been concentrated on the day of publication, with daily news predicting returns for only one to two days \cite{heston2017news,ke2019predicting}. Comparisons with general-purpose LLMs have been inconsistent, with results depending on prompting and fine-tuning \cite{fatouros2023transforming,arxiv2024optimizing}. Few studies combine a multi-source evaluation with a matched ablation across horizons.

\section{Methodology}
\subsection{Data}
News headlines were obtained from FNSPID \cite{dong2024fnspid} and daily prices from Yahoo Finance. One stock with news coverage in every year of 2018--2023 was selected from each of five sectors: INTC, SBUX, WFC, CVX and MRK. This gave 34,810 headlines and 1,509 trading days per stock. Headlines were assigned to trading days, with those on non-trading days moved to the next trading day. As 99.8\% of the timestamps were exactly midnight UTC, a headline dated $D$ could not be separated into before or after the market close of day $D$, and a robustness check was therefore included.

\subsection{Sentiment and Features}
Each headline was scored with a pretrained FinBERT model. Its score was $s = P(\text{positive}) - P(\text{negative})$. Six comparison models were also applied, namely FinBERT-tone, DistilRoBERTa-fin, RoBERTa-Twitter, VADER, BART-MNLI and Qwen2.5-1.5B. Each read one headline in a single step, without tools or multi-step reasoning. Daily sentiment features comprised the mean and standard deviation of the score, the share of positive and negative headlines, the log headline count and a score decayed across days with a half-life of three trading days. A second set was built from headlines naming the company. Fourteen price features covered returns, volatility, moving-average ratios, the Relative Strength Index, volume and daily range.

\subsection{Models and Evaluation}
The forward return over $h$ of 1, 3 and 5 trading days was converted into down, stable and up classes, with the stable band set to the 33rd percentile of absolute returns in the training window. XGBoost \cite{chen2016xgboost} and LSTM \cite{hochreiter1997long} models were trained for each stock and horizon, using fixed hyperparameters shared by all feature sets. As the plain LSTM rarely predicted the down class, a class-weighted LSTM was added. Models were tested on each of 2021, 2022 and 2023 and were trained on all earlier data, with the last $h$ training rows removed.

Macro-F1 was the primary metric, and RMSE and directional accuracy were also recorded. Differences from the price-only model were tested with a paired moving-block bootstrap \cite{kunsch1989jackknife}, using blocks of 10 days and 2,000 resamples, and $p$-values were corrected with the Benjamini--Hochberg procedure \cite{benjamini1995controlling}. The experiments were repeated with sentiment lagged by one trading day. Across outlets, market agreement was defined as the share of non-neutral headlines whose sentiment sign matched the return sign. It was computed from the close of day $T-1$ to $T+1$ (two-day), from $T$ to $T+1$ (forward) and from $T+1$ to $T+2$ (lagged). The Spearman correlation between daily sentiment and returns before, on and after the news day was also measured, together with a trading rule that held the top sentiment quintile long and the bottom quintile short for one day, with cut-offs from 2018--2020 and a cost of 10 basis points per side.

\section{Results}
\subsection{Sentiment Models on Labelled Text}
On the Twitter Financial News dataset \cite{zeroshot2022twitter}, on which none of the models was trained, DistilRoBERTa-fin obtained the highest macro-F1 (0.706), followed by FinBERT (0.668) and FinBERT-tone (0.661). The larger Qwen2.5-1.5B obtained 0.569 and was about 15 times slower than FinBERT, whereas VADER obtained 0.447.

\subsection{Ablation}
All pooled macro-F1 values lay between 0.291 and 0.352, close to the value of about 0.33 expected by chance. At three and five days, the accuracy of XGBoost was below that of always predicting the most common class. No model improved on a zero-return forecast when the stocks were pooled. The effect of sentiment is shown in Table~\ref{tab:paper-gain}.

\begin{table}[H]
\centering
\caption{Change in macro-F1 from adding sentiment, all stocks pooled}
\label{tab:paper-gain}
\begin{tabular}{@{}llrrrr@{}}
\toprule
\textbf{Model} & \textbf{Horizon} & \multicolumn{2}{c}{\textbf{Same-day}} & \multicolumn{2}{c}{\textbf{Lagged}} \\
 & & \textbf{Change} & \textbf{$q$} & \textbf{Change} & \textbf{$q$} \\
\midrule
XGBoost & 1 day & +0.005 & 0.726 & +0.001 & 0.930 \\
 & 3 days & +0.013 & 0.344 & +0.012 & 0.486 \\
 & 5 days & +0.014 & 0.317 & 0.000 & 0.930 \\
\midrule
Class-weighted LSTM & 1 day & 0.000 & 0.992 & $-$0.004 & 0.820 \\
 & 3 days & +0.018 & 0.284 & +0.009 & 0.684 \\
 & 5 days & +0.017 & 0.391 & +0.018 & 0.607 \\
\bottomrule
\end{tabular}
\end{table}

The plain LSTM showed significant gains of +0.025, +0.033 and +0.033 ($q \leq 0.029$). However, the price-only network predicted the down class on only 7--12\% of test days, and sentiment raised this to 20--25\%, while the precision of those predictions stayed close to the share of days that were actually down. Once the class-weighted LSTM removed this imbalance, the gains were no longer significant. For the class-weighted LSTM, sentiment significantly worsened the RMSE at five days, with or without the lag ($q \leq 0.018$).

\subsection{Sentiment Models on Prediction}
No sentiment model produced a significant gain with XGBoost, with same-day (pooled $q \geq 0.648$) or lagged (pooled $q \geq 0.399$) sentiment. FinBERT produced the largest same-day gains at three and five days, but this did not hold after lagging: DistilRoBERTa-fin then led at three days (+0.017), and the 5-day gain of FinBERT fell to zero. The lagged gains at three days were similar for six of the seven models (+0.011 to +0.017), including VADER, which labelled 98.5\% of the headlines as neutral. Neither the headline count nor the sentiment scores alone produced a significant gain ($q \geq 0.389$).

\subsection{Timing and Economic Value}
The correlation of daily sentiment with the same day's return was 0.133 (95\% CI 0.110--0.160) and was significant for all five stocks. It was 0.046 on the previous day, 0.027 on the next day and 0.004 on the day after, the same ordering as reported by \cite{ke2019predicting}. In 2021--2023, the trading rule earned 24.7\% a year before costs when holding on the next day, which could be traded only if the news preceded the close, and 11.8\% (not significant) on the second day, which could always be traded. After costs, both lost money ($-$20.8\% and $-$33.7\% a year).

\subsection{News Sources}
Outlets other than Nasdaq were available only until June 2020, so this comparison was restricted to 2018--2020. Market agreement is shown in Table~\ref{tab:paper-source}. Daily sentiment from Nasdaq and Benzinga correlated at only 0.13 across 1,051 stock-days. Separate XGBoost models using each outlet, tested from January to June 2020, showed no significant gain.

\begin{table}[H]
\centering
\caption{Market agreement of FinBERT sentiment (\%)}
\label{tab:paper-source}
\begin{tabular}{@{}llrrrr@{}}
\toprule
\textbf{Stocks} & \textbf{Outlet} & \textbf{Non-neutral} & \textbf{Two-day} & \textbf{Forward} & \textbf{Lagged} \\
\midrule
INTC & GuruFocus & 74 & 71.6 & 55.4 & 51.4 \\
 & Nasdaq & 1,077 & 56.7 & 48.5 & 49.3 \\
 & Zacks & 246 & 53.3 & 49.6 & 50.8 \\
\midrule
Four stocks & Benzinga & 1,474 & 62.1 & 51.9 & 49.9 \\
 & Nasdaq & 2,713 & 54.8 & 50.5 & 50.4 \\
\bottomrule
\end{tabular}
\end{table}

\section{Discussion}
FinBERT sentiment did not significantly improve the prediction of stock movement beyond price and volume features, and this was unchanged when the sentiment was lagged to remove possible look-ahead. These results are consistent with the efficient market hypothesis \cite{fama1970efficient}. The news-source results point in the same direction, as the 62\% and 55\% agreement of Benzinga and Nasdaq fell to about 50\% when the publication day was excluded. Sentiment appeared to describe price moves as they occurred, rather than to anticipate later ones, and the remaining next-day effect was too small to trade after costs.

These results do not prove that sentiment has no predictive value, as the confidence intervals of about $\pm 0.013$ to $\pm 0.016$ could conceal a small effect. They are weaker than some results in the literature \cite{gu2024predicting,dong2024fnspid}, which may reflect the matched baselines, walk-forward testing and corrected tests used here, although differences in stocks, periods and models prevent a direct comparison. The plain LSTM also showed that a significant gain in macro-F1 can arise from a change in how often a class is predicted. The main limitations were date-only timestamps, the use of headlines only, the dominance of Nasdaq in the data (85.7\% of headlines), and the selection of five stocks with high news coverage.

\section{Conclusion}
A matched ablation found no reliable gain from FinBERT sentiment over price features at horizons of one to five trading days, and the most accurate sentiment model on labelled text could not be shown to be the most useful for prediction. Future research should use news with exact publication times, a longer test period to detect smaller effects, and smaller stocks, for which prices may absorb news more slowly.

\section{Acknowledgment}
The author thanks Ts. Dr. Nurul Aida Binti Osman for supervising this project, and the authors of FNSPID and of the open sentiment models for making their data and models publicly available.

% Prints every entry cited with \cite{...}, numbered in citation order.
\section{References}
{\singlespacing
\printbibliography[heading=none]
\par}

\end{document}
